
\def\PUDcodigo{ACE001}

\if\COMPILAR0
   \def\PUDcodacad{04507.51}
\else
   \def\PUDcodacad{04506.51}
\fi

\def\PUDdisciplina{Libras}

\def\PUDsemestre{Opt}

\def\PUDhoras{80}

%NOVOS ---------------------------------------------------
\def\PUDhorasTeoria{80}
\def\PUDhorasPratica{0}

\def\PUDinterECA{---}

\def\PUDatuaisECA{---}

\def\PUDrecursos{Projetor multimídia; Quadro branco e pincel.}

%----------------------------------------------------------

\def\PUDcreditos{4}

\def\PUDprerequisito{---}

\def\PUDpreacad{---}

\def\PUDobjetivos{Conhecer as características fundamentais da Língua Brasileira de Sinais (Libras), numa proposta comunicativa básica que permita a compreensão e conversação em situações cotidianas.}

\def\PUDementa{Introdução: aspectos clínicos, educacionais e sócio antropológicos da surdez. A Língua de Sinais Brasileira - Libras: características básicas da fonologia. Noções básicas de léxico, de morfologia e de sintaxe com apoio de recursos audiovisuais; Noções de variação.}

\def\PUDconteudo{ &  A Língua de Sinais Brasileira e a constituição linguística do sujeito surdo - Breve introdução aos aspectos clínicos, educacionais e sócio-antropológicos da surdez; Nomeação de pessoas e de lugares em Libras; Noções gerais da gramática de Libras; Prática introdutória de Libras: alfabeto manual ou datilológico; 
\\ 
 &  Noções básicas de fonologia e morfologia da Libras; Parâmetros primários da Libras; Parâmetros secundários da Libras; Componentes não-manuais; Aspectos morfológicos da Libras: gênero, número e quantificação, grau, pessoa, tempo e aspecto;
Prática introdutória de Libras: diálogo e conversação com frases simples; 
\\ 
 &  Noções básicas de morfossintaxe; A sintaxe e incorporação de funções gramaticais; O aspecto sintático: a estrutura gramatical do léxico em Libras; Verbos direcionais ou flexionados; A negação em Libras; Prática introdutória de Libras: diálogo e conversação com frases simples.
\\
 &  Noções básicas de variação; Características da língua, seu uso e variações regionais; A norma, o erro e o conceito de variação; Tipos de variação linguística em Libras; Prática introdutória de Libras: registro videográfico de sinais. }

\def\PUDmetodo{Aulas teóricas; exibição de vídeos; expressão gestual e corporal. Os seguintes recursos poderão ser utilizados: Quadro e pinceis; Projetor de Multimídia e material impresso. A Prática de Componente Curricular de Ensino poderá ser ministrada através de: aulas expositivas, criação e aplicação de técnicas de ensino, apresentação de seminários, elaboração de material didático e realização de projetos em instituições com surdos.}

\def\PUDavalia{A avaliação será desenvolvida ao longo do semestre, de forma processual e contínua, a partir da produção de diálogos em Libras, contação de histórias em Libras, produção de relatos em Libras, produção de objetos de aprendizagem (vídeos, materiais didáticos etc) em Libras, realização de projetos de intervenção em instituições com surdos e participação nas atividades propostas. Serão utilizados para a avaliação os seguintes critérios: conhecimento individual sobre temas relativos aos assuntos estudados em sala; grau de participação do aluno em atividades que exijam produção individual e em equipe; planejamento, organização, coerência de ideias e clareza na elaboração de trabalhos escritos ou destinados à demonstração do domínio dos conhecimentos técnico-pedagógicos e científicos adquiridos; criatividade e o uso de recursos diversificados; domínio de atuação discente (postura e desempenho). 

Os aspectos quantitativos da avaliação ocorrerão de acordo com o Regulamento da Organização Didática (ROD) do IFCE. \vspace{12pt}}

\def\PUDbibbasica{ &  \href{www.biblioteca.ifce.edu.br/index.asp?codigo_sophia=24272}{Quadros}, Ronice Muller.  Língua de sinais brasileira: estudos linguisticos.  Volume único.  Porto Alegre, 2004.  Editora Artmed.  ISBN: 8536303085. 
\\
&   \href{www.biblioteca.ifce.edu.br/index.asp?codigo_sophia=62419}{Gesser}, Andrei.  Libras? Que língua é essa?: crenças e preconceitos em torno da língua de sinais e da realidade surda.  São Paulo: Parábola, 2009.  ISBN: 9788579340017.
%&  Sacks, Oliver W.  Vendo vozes: uma viagem ao mundo dos surdos.  Volume único.  Editora Companhia das Letras.  São Paulo, 2010.  ISBN: 8571647798.  
\\ 
 &   \href{www.biblioteca.ifce.edu.br/index.asp?codigo_sophia=62784}{Capovilla}, Fernando Cesar.  Enciclopédia da língua de sinais brasileira 2: o mundo do surdo em libras: artes e cultura, esportes e lazer.  Volume 2.  Editora EDUSP.  ISBN: 9788531408496. }

\def\PUDbibcomplementar{ &  \href{www.biblioteca.ifce.edu.br/index.asp?codigo_sophia=62503}{Honora}, Márcia.  Livro ilustrado de língua brasileira de sinais: desvendando a comunicação usada pelas pessoas com surdez.  São Paulo, SP: Ciranda Cultural, 2010.  352p.  ISBN: 9788538014218.
\\
&  \href{www.biblioteca.ifce.edu.br/index.asp?codigo_sophia=56532}{Pereira}, Maria Cristina da Cunha.  Libras: conhecimento além dos sinais.  E-book.  ISBN: 9788576058786. 
\\ 
 &  \href{www.biblioteca.ifce.edu.br/index.asp?codigo_sophia=58860}{Silva}, Rafael Dias. Língua brasileira de sinais libras.  E-book. ISBN: 9788543016733.
\\
 &  \href{www.biblioteca.ifce.edu.br/index.asp?codigo_sophia=57213}{Smith}, Adam.  A mão invisível.  E-book.  Cia. das Letras.  132p.  ISBN: 9788563560698.
\\
 &  \href{www.biblioteca.ifce.edu.br/index.asp?codigo_sophia=23955}{Sacks}, Oliver W.  Vendo vozes: uma viagem ao mundo dos surdos.  São Paulo, SP: Companhia das Letras, 2010.  215p.  ISBN: 9788535916089. }

\def\PUDautor{Juliana Brito}

\def\PUDdata{2017-08-11}

\def\PUDrevisao{2}

\def\PUDrevisor{Samuel Vieira Dias}

\def\PUDdatarevisao{2017-08-11}

\PUDcorpo
